\documentclass[conference]{IEEEtran}
\usepackage[UTF8]{ctex}
\usepackage{amsmath,amssymb}
\usepackage{graphicx}
\usepackage{hyperref}
\usepackage{booktabs}
\usepackage{algorithm}
\usepackage{algpseudocode}
\usepackage{cite}

\title{A Context-Aware Multi-Agent Framework for Academic Writing Assistance}
\author{
  \IEEEauthorblockN{First Author}
  \IEEEauthorblockA{Department of Computer Science\\
  University Name, City, Country\\
  Email: author@example.com}
}

\begin{document}

\maketitle

\begin{abstract}
Academic writing is a cognitively demanding task that requires simultaneous attention to argument structure, citation placement, language quality, and domain-specific conventions. This paper presents AutoResearch, a context-aware multi-agent framework that decomposes academic writing into modular subtasks handled by specialized agents. The framework employs a context engine that builds on-demand context packages from parsed LaTeX documents, enabling agents to operate with precise, localized understanding of the manuscript. We evaluate the framework through a multi-stage pipeline covering draft generation, language polishing, and consistency checking. Experimental results demonstrate improved coherence and reduced manual revision effort compared to single-pass generation approaches.
\end{abstract}

\begin{IEEEkeywords}
academic writing, multi-agent systems, context-aware computing, LaTeX, document understanding
\end{IEEEkeywords}

\bibliographystyle{IEEEtran}
\bibliography{refs}

\end{document}
